\documentclass[10pt,oneside]{book}

\usepackage{tin_tome}
\usepackage{amsmath}
\usepackage{amsfonts}
\usepackage{mathtools}
\usepackage{enumitem}
\usepackage{needspace}

\setcounter{secnumdepth}{0}
\setlength{\parindent}{0pt}
\setlength{\parskip}{0.26em}
\setstretch{1.03}
\raggedbottom
\sloppy

\titlespacing*{\section}{0pt}{0.6em}{0.45em}
\titlespacing*{\subsection}{0pt}{0.6em}{0.25em}

\newcommand{\lean}[1]{{\footnotesize\color{muted}\texttt{\detokenize{#1}}}}
\newcommand{\statement}[3]{%
  \Needspace{3\baselineskip}%
  \noindent{\color{chapterbrown}\bfseries #1}\par
  \lean{#2}\par
  #3\par\vspace{0.25em}%
}
\newcommand{\support}[1]{%
  {\footnotesize\color{muted}\raggedright #1\par}\vspace{0.18em}%
}
\newcommand{\problem}[2]{%
  \clearpage
  \section*{#1. #2}
  \markright{#1. #2}
}

\newcommand{\RR}{\mathbb{R}}
\newcommand{\NN}{\mathbb{N}}
\newcommand{\Prob}{\mathbb{P}}
\newcommand{\DR}{\operatorname{DR}}
\newcommand{\ST}{S_{\mathrm{T}}}
\newcommand{\etaL}{\eta_{\mathrm{lyap}}}
\newcommand{\PhiR}{\Phi}
\newcommand{\Pois}{\operatorname{Pois}}
\newcommand{\Binom}{\operatorname{Binom}}
\newcommand{\BigO}{\mathcal{O}}
\newcommand{\supp}{\operatorname{supp}}
\newcommand{\univ}{\mathrm{univ}}
\newcommand{\Icc}{\operatorname{Icc}}
\newcommand{\atTop}{\mathrm{atTop}}

\begin{document}
\thispagestyle{plain}

\begin{center}
{\Large\bfseries\color{chapterbrown} TINProofs C1-C5}\par
{\large Formal Lemmas and Theorems}\par
\vspace{0.4em}
{\small J. Councilman}\par
{\small May 24, 2026}
\end{center}

\vspace{0.4em}
{\color{rulecolor}\hrule height 0.5pt}
\vspace{0.8em}

\section*{C1. Commodity Hull Theorem}
\markright{C1. Commodity Hull Theorem}

\textsc{Notation.}
For a feasible path \(p=(T_p,Q_p)\), define
\[
  A_p(\lambda)=-\log Q_p+\lambda T_p,
  \qquad
  V_p(\lambda)=\log Q_p-\lambda T_p.
\]
For a finite nonempty hull \(C\),
\[
  F_C(\lambda)=\inf_{p\in C}A_p(\lambda),
  \qquad
  V_C(\lambda)=\sup_{p\in C}V_p(\lambda).
\]

\statement{Affine path action and value}
{action_affine; value_affine; action_neg_value}
{For \(a,b\geq 0\) with \(a+b=1\),
\[
  A_p(a\lambda_1+b\lambda_2)
  =aA_p(\lambda_1)+bA_p(\lambda_2),
  \qquad
  V_p(a\lambda_1+b\lambda_2)
  =aV_p(\lambda_1)+bV_p(\lambda_2),
\]
and \(A_p(\lambda)=-V_p(\lambda)\).}
\support{Convexity and concavity corollaries: \lean{action_convexOn; action_concaveOn; value_convexOn; value_concaveOn}.}

\statement{Finite hull support and value duality}
{support_concaveOn; valueFn_eq_neg_support; valueFn_convexOn; legendre_duality}
{The finite lower envelope \(F_C\) is concave on \(\univ\), the finite upper envelope \(V_C\) is convex on \(\univ\), and
\[
  V_C(\lambda)=-F_C(\lambda).
\]
This is the formal Legendre-duality surface for the finite path hull.}

\statement{Crossover equality}
{action_eq_at_crossover}
{For distinct exposures \(T_i\neq T_j\), set
\[
  \lambda_{ij}
  =\frac{\log Q_j-\log Q_i}{T_j-T_i}.
\]
Then the path actions agree exactly:
\[
  A_i(\lambda_{ij})=A_j(\lambda_{ij}).
\]}

\statement{Crossover ordering}
{action_lt_below_crossover; action_gt_above_crossover}
{If \(T_i<T_j\), then
\[
  \lambda<\lambda_{ij}\Rightarrow A_j(\lambda)<A_i(\lambda),
  \qquad
  \lambda_{ij}<\lambda\Rightarrow A_i(\lambda)<A_j(\lambda).
\]
Thus the theorem identifies both the crossing point and the side on which each path is selected.}

\problem{C2}{Stochastic Certification Asymmetry}

\textsc{Notation.}
Let \(k>0\) be the reserve-cycle count, let tier thresholds be
\(\tau_i\in(0,1)\), and define the integer cutoff
\[
  x_i=\left\lceil k\tau_i\right\rceil.
\]
For \(X\sim\Binom(n,p)\),
\[
  B_{n,p}(j)=\binom{n}{j}p^j(1-p)^{n-j}.
\]

\statement{Absorbing lowest-tier cutoff}
{ceil_mean_lt_xMin_of_lt_sub_one; median_lt_xMin_of_median_le_ceil_mean; absorbing_lowest_tier_median_lt_cutoff}
{If \(p\geq0\) and
\[
  np < k\tau_i-1,
\]
then
\[
  \left\lceil np\right\rceil < x_i.
\]
Consequently, any median \(m\) with \(m\leq\lceil np\rceil\) lies below the same cutoff.  The lowest-tier theorem is this statement specialized to the final threshold \(i=i_{\min}\).}
\support{The mode consequence is recorded separately as \lean{mode_lt_xMin_of_mode_le_median}.}

\statement{Tier asymmetry as finite sums}
{TierSystem.xMin_pos; binomialCDF_pred_eq_sum_range; tierAsymmetryRatio_eq_binomial_sums}
{The tier-specific downgrade-to-upgrade asymmetry ratio is exactly
\[
  R_{T,T^-}
  =
  \frac{\sum_{j=0}^{x_T-1} B_{n,p}(j)}
       {\sum_{j=x_{T^-}}^{n} B_{n,p}(j)}.
\]
The proof first shows each cutoff \(x_i\) is positive, so the lower tail is the range \(0,\ldots,x_T-1\).}

\statement{Binomial-to-Poisson tail limit}
{binomialPMF_tendsto_poissonPMF; binomialLowerTail_tendsto_poissonLowerTail; binomialUpperTail_tendsto_poissonUpperTail}
{If \(n p_n\to r\), then for each fixed \(k\),
\[
  B_{n,p_n}(k)\longrightarrow e^{-r}\frac{r^k}{k!}.
\]
The corresponding finite lower tails and upper tails converge to the Poisson tails
\[
  L_m(\mu)=\sum_{j=0}^{m}e^{-\mu}\frac{\mu^j}{j!},
  \qquad
  U_k(\mu)=\sum_{j\geq k}e^{-\mu}\frac{\mu^j}{j!}.
\]}
\support{Supporting normalizations and complement identities: \lean{binomialPMF_sum_range_eq_one; poissonUpperTail_eq_one_sub_lower; hasSum_poissonPMF}.}

\statement{Poisson crossover existence}
{poissonCrossoverCondition_iff; poissonCrossover_exists_of_bracket; poissonCrossover_exists_of_reverse_at}
{The limiting crossover equation is
\[
  L_{k-2}(\mu)=U_k(\mu).
\]
If \(a\leq b\), \(U_k(a)\leq L_{k-2}(a)\), and \(L_{k-2}(b)\leq U_k(b)\), then there exists
\[
  \mu\in[a,b]\quad\text{with}\quad L_{k-2}(\mu)=U_k(\mu).
\]
The specialized theorem uses \(a=0\) and any nonnegative right endpoint where the inequality has reversed.}

\statement{Asymptotic delivery-ratio expansion}
{adellJodraAsymptoticInput; drStar_expansion_of_lambda_expansion}
{Under the named asymptotic input
\[
  \lambda_k^*=k-\frac{5}{6}+\BigO(k^{-1}),
\]
the selected delivery ratio
\[
  \DR_k^*=\frac{\lambda_k^*}{k}
\]
satisfies
\[
  \DR_k^*
  =
  1-\frac{5}{6k}+\BigO(k^{-2}).
\]
This theorem is conditional on the external Poisson-median asymptotic boundary.}

\statement{Gap tier}
{no_achievable_in_gap; gap_tier}
{If two adjacent ceiling cutoffs coincide,
\[
  x_i=x_j,
\]
then no integer \(X\) satisfies
\[
  x_i\leq X<x_j.
\]
}

\problem{C3}{Lyapunov Radius of Validity}

\textsc{Notation.}
The setup stores positive constants
\(\alpha_{\min},\alpha_{\max},q,r_0,\lVert A\rVert,M\) and defines
\[
  r_*=\min\left(r_0,\frac{q}{2\lVert A\rVert M}\right),
  \qquad
  c_*=\frac{q}{2\alpha_{\max}}.
\]
For a perturbation point \(P\), write \(\rho=\lVert\delta\rVert\), \(V=P.V\), and \(\dot V=P.Vdot\).

\statement{Eigenvalue sandwich}
{eigen_sandwich_lower; eigen_sandwich_upper; eigen_sandwich}
{The Lyapunov quadratic value satisfies
\[
  \alpha_{\min}\rho^2\leq V\leq\alpha_{\max}\rho^2.
\]
The lower and upper inequalities are available both separately and as a conjunction.}

\statement{Derivative decomposition and pointwise bounds}
{vdot_decomposition; quadratic_term_lower; negative_quadratic_term_upper; remainder_abs_bound; remainder_upper_bound}
{The derivative decomposes as
\[
  \dot V=-Q(\delta)+R(\delta),
\]
with
\[
  q\rho^2\leq Q(\delta),
  \qquad
  |R(\delta)|\leq \lVert A\rVert M\rho^3.
\]
Hence \(-Q(\delta)\leq -q\rho^2\) and \(R(\delta)\leq \lVert A\rVert M\rho^3\).}

\statement{Critical-ball derivative bound}
{vdot_le_quadratic_plus_cubic; vdot_bound_on_ball}
{Combining the previous inequalities gives
\[
  \dot V\leq -q\rho^2+\lVert A\rVert M\rho^3.
\]
Inside the critical ball \(\rho<r_*\), the cubic term consumes at most half the quadratic margin:
\[
  \dot V\leq -\frac{q}{2}\rho^2.
\]}

\statement{Main radius theorem}
{criticalRadius_pos; lyapunov_radius_of_validity}
{The radius \(r_*\) is positive.  If \(\rho<r_*\), then
\[
  \dot V\leq -c_* V
  =
  -\frac{q}{2\alpha_{\max}}\,V.
\]
This is the formal Lyapunov radius-of-validity theorem.}

\statement{Sublevel-set corollaries}
{sublevel_inside_critical_ball; sublevel_boundary_vdot_nonpos}
{If \(V\leq V_0\) and
\[
  V_0<\alpha_{\min}r_*^2,
\]
then \(\rho<r_*\).  On the boundary \(V=V_0\) of such a sublevel set,
\[
  \dot V\leq0.
\]}

\problem{C4}{Temporal Transport Factorization}

\textsc{Notation.}
Let \(D\subseteq F\) be delivery and feasibility events in a probability space.  Define
\[
  \DR=\Prob(D),
  \qquad
  \ST=\Prob(F),
  \qquad
  \eta=\Prob(D)/\Prob(F).
\]

\statement{Delivery implies feasibility}
{delivery_inter_infeasible_eq_empty; delivery_zero_when_infeasible; delivery_measure_le_feasible}
{The event-level inclusion \(D\subseteq F\) yields
\[
  D\cap F^c=\emptyset,
  \qquad
  \Prob(D\cap F^c)=0,
  \qquad
  \Prob(D)\leq\Prob(F).
\]}

\statement{Exact factorization}
{exact_factorization; exact_factorization_real}
{Whenever \(\ST\neq0\),
\[
  \DR=\ST\,\eta.
\]
The theorem is proved both in the extended nonnegative-real measure layer and in the real-valued ratio layer.}

\statement{Augmentation monotonicity}
{st_monotone; stReal_monotone}
{For an augmentation \(F\subseteq F'\) with the same underlying measure,
\[
  \ST\leq\ST'.
\]
The same monotonicity holds after conversion to real-valued factors.}

\statement{Braess localization}
{braess_in_eta; braess_in_etaReal}
{If an augmentation has \(\ST'>0\), preserves the measure, and nevertheless decreases delivery ratio,
\[
  \DR'<\DR,
\]
then the decrease is localized in the efficiency factor:
\[
  \eta'<\eta.
\]}

\problem{C5}{Three-Factor Sparse Law}

\textsc{Notation.}
The scalar setup records \(\ST\), \(\eta\), \(E[H]\), a Lyapunov exponent \(\lambda<0\), and \(\DR=\ST\eta\).  Define
\[
  \etaL=\exp(E[H]\lambda),
  \qquad
  \PhiR=\eta/\etaL.
\]

\statement{Exact three-factor law}
{three_factor}
{The sparse-law delivery ratio factors as
\[
  \DR=\ST\,\etaL\,\PhiR.
\]
This is an algebraic refinement of \(\DR=\ST\eta\) by inserting the chain attenuation factor and residual routing distortion.}

\statement{Chain attenuation properties}
{etaLyap_pos; etaLyap_lt_one; phi_pos}
{The chain attenuation is strictly positive,
\[
  0<\etaL,
\]
and, when \(E[H]>0\) and \(\lambda<0\),
\[
  \etaL<1.
\]
If \(\eta>0\), then \(\PhiR>0\).}

\statement{Negative mean log-hop exponent}
{lyapunov_exponent_neg}
{For finitely many hop probabilities \(p_i\in(0,1)\) with \(n>0\),
\[
  \frac{1}{n}\sum_{i=1}^{n}\log p_i < 0.
\]
This supplies the sign condition used by the chain attenuation factor.}

\statement{Morphology vocabulary}
{morphologySlope; isTrap; isCluster}
{For two residual measurements \((\Phi_1,E_1)\) and \((\Phi_2,E_2)\) with positive residuals and \(E_1\neq E_2\), define
\[
  \gamma
  =
  \frac{\log \Phi_2-\log \Phi_1}{E_2-E_1}.
\]
The formal classes are \(\gamma<0\) for a trap and \(\gamma>0\) for a cluster.}

\end{document}
